% common/user-blocks.tex
% Shared sections of the user part: front-matter blocks and bibliography.
% Uses \unishort, \uniotherlangnote and \unibibdefault, set by the guide.
\section{Abstract, originality, acknowledgements, dedication}
\label{sec:blocks}

The core class provides four front-matter blocks. They are written the
same way at every university; the \unishort{} package only sets their look.

\begin{tabularx}{\textwidth}{@{}l l Y@{}}
\toprule
\textbf{Block} & \textbf{Options} & \textbf{Default heading (English)} \\
\midrule
\cs{begin}\opt{\{abstractthesis\}} & \opt{language}, \opt{keywords}, \opt{title} & Abstract \\
\cs{begin}\opt{\{originality\}} & \opt{language}, \opt{title} & Statement of Originality \\
\cs{begin}\opt{\{acknowledgements\}} & \opt{language}, \opt{title} & Acknowledgments \\
\cs{dedication}\opt{[...]\{text\}} & \opt{language} & none: the text alone on its page \\
\bottomrule
\end{tabularx}

\begin{itemize}
  \item \opt{language=} writes the block in another language: its heading,
    the ``Keywords'' label and the hyphenation follow that language.
    Without it, the block uses the main language of the thesis.
    \uniotherlangnote
  \item \opt{keywords=} prints a list of keywords after the abstract.
  \item \opt{title=} replaces the default heading.
\end{itemize}

Abstracts, statements of originality and acknowledgements are listed in
the table of contents. An English abstract followed by an Italian one:

\begin{lstlisting}
\begin{abstractthesis}[keywords = {corporate finance, valuation}]
  The English abstract.
\end{abstractthesis}
\begin{abstractthesis}[language = italian,
                       keywords = {finanza aziendale, valutazione}]
  L'abstract in italiano.
\end{abstractthesis}
\dedication{To my family.}
\end{lstlisting}

\section{Bibliography}
\label{sec:bib}

The class loads \pkg{biblatex} with \pkg{biber} for you; do not load
\pkg{biblatex} again in your preamble. Its style is, in order of priority:
the options you give with the class option \opt{biblatex-options}; the
university's default (\unibibdefault); otherwise a numeric style.

\begin{lstlisting}
\documentclass[..., biblatex-options={style=apa}]{academicthesis}
\addbibresource{references.bib}   % in the preamble
...
\printbibliography                % where the bibliography goes
\end{lstlisting}

If \pkg{biblatex} is not installed, the class warns and uses \pkg{natbib}
with Bib\TeX{} behind the same two commands, so the thesis still compiles.
